\section{Discussion}
\label{sec:discussion}

\subsection{Why the results are what they are}
\label{sec:disc-why}

The results follow from one property of the test, that it reads a change as an event with no size, and from the kind of parameter it reads. A stated parameter changes only when someone decides, so the indicator of a change is a record of decisions, and a coincidence of a decision with a breach is the event the test is built to see. The detection of one, in every cell, is the arithmetic of Proposition~\ref{prop:separation} and not an empirical finding about banks, and its value is that it holds without a model of the parameter. The false flags are the other side of the same arithmetic: a decision that happens to fall on a breach date is indistinguishable from one made because of it, so the false-flag rate is set by how busy the governance calendar is ($\lambda$) and how long the breach lasts ($m$), and it is not a property of the test that can be tuned. A re-estimated parameter changes at every date because the window moves, so the indicator of a change records the arithmetic of the estimator and the coincidence with a breach is the coincidence with a quarter-end. The degeneracy is then complete up to the ties at the cap, which is why the area under the ROC curve of the count against a planted override is \fitAucCEngine{} and does not move with the size of the override.

The tolerance rule works for the reason a tolerance always does: it replaces the question whether the parameter changed by the question whether it changed by more than the estimator changes it. Its power is the power of the maximum of the changes at the breach dates against a null whose maximum is read from a population, and Proposition~\ref{prop:gaussian} shows that this power rises slowly with the size of the override in units of the natural change and falls with the number of breach dates, as the experiments find: at the tolerance of \deltaB{} years the detection is \fitAddCTol{} at 2 years and \fitAddDTol{} at 4. The conformal calibration holds the size on average, and its size at a different breach prevalence depends on the exchangeability that the proposition names; the measured departure, \fitMisCTol{} against 5\,\% at a prevalence of 0.50, is inside the spread of a single calibration. On the real series the tolerance is conservative because the simulated population has changes larger than the real ones; the ranking is perfect at 2 years and the calibrated threshold is too high. Both facts have the same cause and the same remedy, a null population closer to the bank than a block bootstrap of its own rows.

\subsection{What a red light means and how to use the test}
\label{sec:disc-use}

On a stated parameter a red light means that an assumption was changed at a date at which a limit was breached. It does not mean that the change was made because of the breach, and Equation~\eqref{eq:false-flag} says how often it was not: at the rate of unrelated changes of the base cell and a mean of \statedMeanM{} breach dates, \statedNullEngine{} of the histories with no planted change are red. The right use is the one the engine's documentation gives, a prompt for a review of the minutes and the model file at the flagged date, with the false-flag rate read from the closed form for the bank's own governance calendar. A green light on a stated parameter is stronger evidence than a red one: no decision of any size was made at a breach date, which is a fact about the file and not about a distribution.

On a re-estimated parameter the exact test should not be used as a finding. The stored rows of case C14 are red because the weighted average maturity is fitted afresh at each date and the limit is breached at each, and the product's statement that those rows show the test operating on the parameter, and not that the bank altered anything, is the correct reading and follows from Proposition~\ref{prop:degeneracy}. Two uses remain. The first is to apply the test to the parameter that the bank states, the core maturity or the overlay that it adds to the fit, and not to the fit; for a bank that holds an override on top of a fitted value, the history of the override is a stated parameter. The second is the tolerance rule on the fit, calibrated on a null population of banks with the same estimator, length and breach process, and read as a ranking of how large the largest change at a breach date is relative to what the estimator produces; its value is the area under the ROC curve, and a threshold is a choice of a false-flag rate. Neither of these is implemented in the engine, which has the exact test only, and the second needs a population that a single validator may not have. The weight stability test of the engine is the existing check of the fit and it is red in about half of the null banks, so it too says more about the estimator than about the bank.

\subsection{Negative results and their implications}
\label{sec:disc-negative}

Three predictions were refuted and one family of overrides was detected only weakly by every method on the simulated banks and moderately by a test of the level on the real series. The refutation of H5 means that the test cannot be presented as a hypothesis test with a conventional size; it is a screen whose false-flag rate depends on the calendar. The refutation of H6 means that its output on a fitted parameter says little about an override: its area of \fitAucCEngine{} comes from the ties of the estimate at its cap, and it is \realEngineAuc{} on the real fitted series. H8 fails on its point estimate, \fitMisCTol{} (\fitMisCTolLo{} to \fitMisCTolHi{}\,\%) at a prevalence of 0.50 for a tolerance calibrated at 0.25, which is inside the spread of a single calibration; the data show the direction and not a loss of size beyond what the sampling of a calibration produces: a bank whose limit is breached more often needs a larger tolerance (\deltaA{}, \deltaB{} and \deltaC{} years at the three prevalences), since the maximum is taken over more dates. The result for H9 means that the dating convention, by which a change counts only at the date whose window contains the breach, is what keeps an anticipated override from being seen: a persistent override made at the date before a breach begins produces one change, at a date that is not a breach date, which the same-date rule cannot count and the look-ahead variant can, provided the change is larger than what the estimator produces, since the variant is a tolerance rule. The override toward the longest tranche is invisible in about a fifth of the histories because the fitted weights already sit there, and its effect on the others is a small lengthening and a damping of later changes; on the simulated banks no detector that reads the series exceeds an area of \wTopAuc{}, while on the real series, whose changes are steadier, a test of a shift in the level reaches \wRealTopAuc{}. A validator who wants to see such an override can therefore read the level of the series with a break test, which the tests of the changes do not do; whether reading the weights themselves, against the weights that the unconstrained fit would give, would see more of it was not tested here (the stability test of the engine, which reads them, reaches \realWAucStab{} on the real series). These results are not defects of the implementation. They are the properties of a coincidence rule, and the measured rates give a validator the numbers needed to decide how much weight a red or a green light deserves on a given parameter.
